\documentclass[12pt]{article}

\usepackage[T1]{fontenc}
\usepackage[utf8]{inputenc}
\usepackage{lmodern}
\usepackage[ngerman]{babel}
\usepackage[autostyle]{csquotes}
% \usepackage[authoryear,round,<other options>]{natbib} % Required for authoryear citation style
\usepackage[style=authoryear,sorting=nyt, backend=biber]{biblatex}
\addbibresource{references.bib}

\usepackage{framed}



\title{Bachelorarbeit \\ Autonomes Landen von Drohnen in Agrarlandschaften basierend auf Deep-Reinforcement-Learning}
\author{Fritz Körner}

\date{Juni 2026}

\begin{document}

\maketitle

\section{Einleitung}
Der \textit{Klimawandel} ist eine langfristige Veränderung der durchschnittlichen Wetterverhältnisse, die das lokale, regionale und globale Klima der Erde prägen \autocite{WhatClimateChange2022}. Solche Veränderungen können auf natürliche Weise stattfinden, jedoch ist seit dem 19. Jahrhundert die Menschheit der Hauptantrieb für Klimaveränderungen \autocite{nationsWhatClimateChange}, weshalb auch vom menschengemachten Klimawandel gesprochen wird. Der menschengemachte Klimawandel stellt wohlmöglich eine der größten Herausforderungen der Moderne dar. Ausgelöst durch hohe Treibhausgasemissionen gab es bereits in den Jahren 2011-2020 erhöhte Oberflächentemperatur von 1.1°C im Vergleich zum Vor-Industriellen Zeitalter \autocite{calvinIPCC2023Climate2023}. Die Auswirkungen der globalen Erwärmung sind unter anderem vermehrtes Auftreten von Hitzewellen, Starkregen und Dürren \autocite{calvinIPCC2023Climate2023}. Die Auswirkungen des Klimawandels sind auch in der Landwirtschaft sichbar, hierzu zählen unter anderem erhöter Schädlings- und Krankheitsdruck, erhöhter Wasserbedarf \autocite{ClimateChangeTrends}. Es ist auch abzusehen das mit steigenden Temperaturen weltweit Landwirtschaftliche Erträge sinken werden \autocite{hultgrenImpactsClimateChange2025}. Zusätzlich steht die Landwirtschaft vor weiteren Herausforderungen: Eine immer weiter wachsende Weltbevölkerung, Fachkräftemangel und wirtschaftlicher Druck \textbf{QUELLEN}.
Eine Vielversprechende Chance zur Bewältigung dieser Herausforderungen bringt dabei die Forschungsdomände "Digitale Landwirtschaft" mit sich \textbf{QUELLEN}. Innerhalb der digitalen Landwirtschaft sind Drohnen ein viel erforschtes Thema. Neben der digitalen Landwirtschaft sind Drohnen auch in anderen Bereichen der Robotik ein sehr Interessantes Forschungsobjekt. Durch ihre vielseitigen Einsatzmöglichkeiten erweisen sie sich als moderne Problemlöser wie Beispielsweise im SearchAndRescue, für Notfalllieferungen für Medikamente und 3.??? \textbf{HIER QUELLEN UND KONKRETE BEISPIELE}. In der Landwirtschaft können Drohnen zum Ausbringen von Saatgut und Pflanzenschutzmitteln genutzt werden, für verschiedene Remote-Sensing Szenarien oder zum autonomen Ernten von Weintrauben \textbf{QUELLEN, unter anderem in Fue 2020}. Allgemein bieten Drohnen eine erfolgsversprechende Aussicht die Landwirtschaft nachhaltiger, effizienter und resilienter tu gestalten \textbf{ECHTE ZITAT: Unde, Kurkute und Chavan 2025; Guebsi, Mami und Chokmani 2024}. Trotz umfangreicher wissenschaftlicher Arbeit zu Drohnen gibt es noch viele Themengebiete mit Entwicklungspotential. Besonders der Bereich der autonomen Navigation ist aktuell ein spannendes Themenfeld.

\begin{itemize}
    \item viele verschiedene Ansätze, je nach Einsatzbereich, Sensorik, was wiederum Einfluss auf erschwinglichkeit/ sinnhaftigkeit des Einsatzes von drohnen hat
    \item hat zuletzt vom Aufschwung in der KI Forschung profitiert, besonders vision-based navigation hat nochmal von verbesserter CNN Technik profitiert
\end{itemize}

Eine zentrale Rolle in einer autonomen Flugmission spielt das Landen von Drohnen. Für die meisten Flugmissionen ist unfallfreies Landen essentiell. Um unfallfreies Landen sicherzustellen müssen die Flugmanöver in der Landephase präzise und verlässlich sein, gleichzeitig darf der Landeanflug nicht zu viel Zeit in Anspruch nehmen. Aufgrund der durchaus hohen Anforderungen und vielzahl verschiedenster Landeszenarien ist autonomes Landen von Drohnen Thema vieler Forschungsarbeiten.

\newpage
\printbibliography
\end{document}